\documentclass[11pt, a4paper]{article}

\usepackage[T1]{fontenc}
\usepackage[utf8]{inputenc}
\usepackage{tgpagella}
\usepackage{tgcursor}
\usepackage{microtype}
\usepackage[spanish, es-noshorthands]{babel}
\addto\captionsspanish{\renewcommand{\tablename}{Tabla}}

\usepackage{amsmath, amssymb}
\usepackage[table, xcdraw]{xcolor}
\usepackage{booktabs}
\usepackage[most]{tcolorbox}
\usepackage[margin=2.5cm]{geometry}
\usepackage{fancyhdr}
\usepackage{listings}
\usepackage{tikz}
\usetikzlibrary{shapes.geometric, arrows.meta, positioning, calc}

% Configuracion de listings con soporte UTF-8 (literate) para el español
\lstdefinestyle{xmlstyle}{
    language=XML,
    basicstyle=\ttfamily\small,
    backgroundcolor=\color{gray!5},
    frame=single,
    rulecolor=\color{gray!40},
    keywordstyle=\color{blue}\bfseries,
    stringstyle=\color{red!70!black},
    commentstyle=\color{green!60!black},
    showstringspaces=false,
    breaklines=true,
    extendedchars=true,
    literate={á}{{\'a}}1 {é}{{\'e}}1 {í}{{\'i}}1 {ó}{{\'o}}1 {ú}{{\'u}}1 {ñ}{{\~n}}1 {¿}{{?`}}1 {¡}{{!`}}1
}
\lstset{style=xmlstyle}

\pagestyle{fancy}
\fancyhf{}
\setlength{\headheight}{16pt}
% Cabeceras acortadas para evitar el solapamiento (Overfull)
\lhead{\textbf{IMT-342} | Robótica}
\chead{Guía Defensa URDF}
\rhead{Semana 6 - Sesión 2}
\lfoot{Prof: M.Sc. Ing. Bernardo Quiroga Turdera}
\rfoot{Página \thepage}
\renewcommand{\headrulewidth}{0.4pt}
\renewcommand{\footrulewidth}{0.4pt}

\definecolor{ucbblue}{RGB}{0, 51, 102}

\begin{document}

\begin{center}
    \Large\textbf{\textcolor{ucbblue}{Guía de Defensa URDF (Hito 1)}}\\[0.2cm]
    \normalsize \textit{Documento Estudiantil - Preparación para Evaluación[cite: 8]}
\end{center}

\begin{tcolorbox}[colback=white, colframe=ucbblue, title=\textbf{Lo que debe demostrar}]
    La defensa no es una presentación sobre cómo se ve el modelo en RViz. Es una sustentación técnica donde deben probar que el modelo URDF es una representación \textbf{cinemáticamente coherente}.
    \begin{center}
        \textbf{Estructura + Movimiento + Razonamiento + Validación + Defensa Individual[cite: 8]}
    \end{center}
\end{tcolorbox}

\section{Distinción entre Colocación y Eje de Movimiento[cite: 8]}
\begin{center}
\begin{tikzpicture}[scale=1.2, thick]
    % Parent Frame
    \coordinate (P) at (0,0);
    \draw[->, ucbblue] (P) -- (1,0) node[right]{$x_{parent}$};
    \draw[->, ucbblue] (P) -- (0,1) node[above]{$z_{parent}$};
    \filldraw[black] (P) circle (1.5pt) node[below left]{Parent};
    
    % Origin translation
    \coordinate (J) at (3,1.5);
    \draw[-Stealth, dashed, gray] (P) -- (J) node[midway, above left]{origin xyz};
    
    % Joint Frame and Axis
    \draw[->, red] (J) -- ($(J)+(1,0)$) node[right]{$x_{joint}$};
    \draw[->, red] (J) -- ($(J)+(0,1)$) node[above]{$z_{joint}$};
    \draw[->, red, line width=1.5pt] (J) -- ($(J)+(0,1)$) node[midway, right]{axis xyz="0 0 1"};
    \draw[-Stealth, red] ($(J)+(0.3,0.5)$) arc (30:150:0.4);
    
    \filldraw[black] (J) circle (1.5pt) node[below right]{Joint};
\end{tikzpicture}
\end{center}
\begin{itemize}
    \item \textbf{\texttt{origin}:} Responde a \textit{¿Dónde está situado el frame de la articulación respecto al parent?}[cite: 8]
    \item \textbf{\texttt{axis}:} Responde a \textit{¿Alrededor de qué vector local se mueve la articulación?}[cite: 8]
\end{itemize}

\section{D-H versus URDF[cite: 8]}
No confundan los marcos D-H con los marcos URDF. En D-H, la colocación sigue las reglas de normales comunes. En URDF, se separan la colocación estática (\texttt{origin}) y el movimiento (\texttt{axis}). Dos modelos pueden diferir en sus marcos intermedios y seguir siendo cinemáticamente equivalentes si su comportamiento final concuerda.

\section{Árbol de Decisión de Troubleshooting (Depuración)[cite: 8]}
Utilice la siguiente lógica si detecta una anomalía en su validación:
\begin{lstlisting}[style=xmlstyle, backgroundcolor=\color{white}]
¿Parsea el archivo XML?
 |- No -> Revisar sintaxis, cierre de etiquetas, comillas.
 |- Si -> ¿El arbol Parent/Child es unicamente serial?
           |- No -> Revisar declaraciones <parent> y <child>.
           |- Si -> ¿El eje gira sobre una direccion errada? -> <axis>
                    ¿El eje rota bien pero desde otro pivote? -> <origin>
                    ¿Todo rota bien pero la malla visual esta salida? -> <visual><origin>
\end{lstlisting}

\section{Secuencia de la Defensa[cite: 8]}
Para cada grupo, el docente solicitará una demostración ágil:
\begin{enumerate}
    \item Mostrar modelo y árbol.
    \item Identificar parent/child de una junta elegida por el profesor.
    \item Explicar su \texttt{origin} y \texttt{axis}.
    \item Ejecutar pruebas de movimiento de juntas aisladas.
    \item Mostrar una configuración de validación con juntas en posiciones diferentes de cero.
    \item Responder preguntas individuales rotativas.
\end{enumerate}

\end{document}
